\documentclass[11pt, letterpaper]{article}

\usepackage[margin=1in]{geometry}
\usepackage{fontspec}
\usepackage{hyperref}
\usepackage{parskip}
\usepackage{microtype}

\setmainfont{Times New Roman}
\hypersetup{colorlinks=true, urlcolor=blue, linkcolor=blue}

\pagestyle{empty}

\begin{document}

% ── Header ──────────────────────────────────────────────────────────────────
\begin{center}
  {\LARGE \textbf{Cover Letter for GSI Position (EAS 531)}} \\[4pt]
  \href{mailto:jiayuanj@umich.edu}{jiayuanj@umich.edu} \quad
  Jiayuan Jiang
\end{center}

\vspace{8pt}

% ── Date & Recipient ────────────────────────────────────────────────────────
\noindent April 6, 2026

\vspace{8pt}

\noindent Prof.\ Kim Diver \\
School for Environment and Sustainability \\
University of Michigan, Ann Arbor

\vspace{12pt}

% ── Salutation ──────────────────────────────────────────────────────────────
\noindent Dear Prof.\ Diver,

\vspace{8pt}

% ── Body ────────────────────────────────────────────────────────────────────

I am writing to express my strong interest in the GSI position for EAS 531: Principles of GIS. I am currently a Master’s student in Geospatial Data Science at the University of Michigan, and my long-term goal is to pursue research in GeoAI, focusing on the integration of spatial data and machine learning for understanding complex geographic processes.

I completed my undergraduate degree in Geographic Information Science at Huazhong Agricultural University, where I received comprehensive training in GIS and remote sensing. My coursework includes Principles of GIS, GPS Technology, Introduction to Remote Sensing, Hyperspectral Remote Sensing, Quantitative Remote Sensing, Active Remote Sensing, UAV Remote Sensing, and GIS Development. This background has provided me with a solid theoretical foundation and strong practical skills in geospatial analysis.

I have extensive experience using ArcGIS Pro and am highly proficient in ArcPy for spatial data processing, including batch raster operations, data preprocessing, and geoprocessing workflows. Through both coursework and research, I have developed strong debugging skills and am familiar with common issues in GIS workflows. I also have experience with QGIS, including plugin-level development, which allows me to understand GIS tools beyond surface-level usage.

I completed EAS 531 in the previous semester and received an A+ grade. This experience allows me to understand the course structure, lab workflow, and common challenges students encounter. I am confident in my ability to guide students through lab exercises, assist with troubleshooting, and provide clear explanations of both concepts and technical steps.

In addition, I am currently working as a Research Assistant under Professor Kai Zhu on a GeoAI project focused on forest disturbance mapping using multi-temporal satellite imagery and machine learning. This experience has strengthened my ability to work with large-scale geospatial datasets, develop analytical workflows, and apply computational methods to real-world environmental problems.

I am particularly interested in this position because I enjoy helping others solve technical problems and learn complex tools. I am comfortable supporting students during lab sessions, assisting with debugging in real time, and ensuring they can complete assignments effectively. I believe my technical background, familiarity with the course, and interest in teaching would allow me to contribute positively to the learning environment.

Thank you for your time and consideration.

\vspace{12pt}
\noindent Sincerely,

\vspace{24pt}
\noindent \textbf{Jiayuan Jiang}

\end{document}